% Source-derived chapter 13: Examples
% Original source: rigid-local-systems-multiplicative-eigenvalue-problem
\chapter{Examples}
\label{rigid-local-systems-multiplicative-eigenvalue-problem-chapter-13}
\source{rigid-local-systems-multiplicative-eigenvalue-problem}

\begin{remark}[Source section]
\label{rigid-local-systems-multiplicative-eigenvalue-problem-chapter-13-source}
\source{rigid-local-systems-multiplicative-eigenvalue-problem}
\source{rigid-local-systems-multiplicative-eigenvalue-problem}
This chapter reproduces the corresponding section of the retrieved
LaTeX source, including its definitions, statements, examples, and
proofs. It has not yet been translated into Lean.
\notready
\end{remark}

% The source section title was: Examples
\label{somemore}
\source{rigid-local-systems-multiplicative-eigenvalue-problem}
In the following we have computed Gromov-Witten numbers in the following way: Convert them into ranks of conformal block bundles by cyclically shifting of the weights $d$ times using the formulas given in Section \ref{QCB}. The ranks of conformal blocks were computed using Swinarski's software package \cite{Swin}.

\subsection{An example from $\op{Gr}(2,4)$}\label{oldie}
\source{rigid-local-systems-multiplicative-eigenvalue-problem}
Let $I^1=\{1,4\}$,  $I^2=I^3=\{1,3\}$, and $d=1$. The Gromov-Witten number $\langle \sigma_{I^1},\sigma_{I^2},\sigma_{I^3}\rangle_1$ equals one by using formulas from Section \ref{stuffed}. Consider the pair $(1,4)$, and the divisor $D(1,4)$ this corresponds to $J^1=J^2=J^3=\{1,3\}$, and $d'=1$. Writing $\mathcal{O}(D(a,j))=\mathcal{B}(\vec{\lambda},\ell)$, we find formulas for $\lambda^1=\lambda^2=\lambda^3$ and $\ell$ using Theorem \ref{brahms}. The formula for $\ell$ is  $\langle \sigma_{J^1},\sigma_{J^2},\sigma_{J^3},\sigma_{J^4}\rangle_2$ with $J^4=\{1,3\}$. Using formulas from Section \ref{stuffed} again, we get
$\langle \sigma_{J^1},\sigma_{J^2},\sigma_{J^3},\sigma_{J^4}\rangle_2=\langle \sigma_{A},\sigma_{A},\sigma_{A},\sigma_{A}\rangle_{2,4}$ which in turns equals
$\langle \sigma_{A},\sigma_{A},\sigma_{A},\sigma_{A}\rangle_{0,0}=2$, where $A=\{2,4\}$, the classical count of the number of lines passing through $4$ general lines in space.

To compute $\lambda^1$ we will need $\langle \sigma_{C},\sigma_{J^2},\sigma_{J^3},\rangle_1$ and  $\langle \sigma_{D},\sigma_{J^2},\sigma_{J^3}\rangle_1$ where $C=\{2,3\}$ and  $D=\{1,4\}$. The latter number is already computed as one. The former is also easily computed to be one by the formulas in Remark \ref{stuffed}.
Therefore $\lambda^1=\lambda^2=\lambda^3= \omega_1+\omega_3=(2,1,1,0)$. The vertex in $P_4(3)$ is $(\vec{a},\vec{a},\vec{a})$ with
$\vec{a} =\frac{(1,0,0,-1)}{\ell}=(\frac{1}{2},0,0,-\frac{1}{2})$. This vertex was first found in \cite[Section 7]{BLocal}. See Section \ref{smalln} for the rank $2$ rigid unitary local system which corresponds to this vertex.



\subsection{An example of Thaddeus}
Consider the ``non-abelian vertex" on page 25 of \cite{thaddy}. Let $\sigma_1,\sigma_2$ and $\sigma_3$ be as in \eqref{dfnsigma}.
The $8\times 8$ matrices
\begin{itemize}
\item $A=\op{diag}(\sigma_1,\sqrt{-1}, -\symbo, \symbo,-\symbo,-1,-1)$
\item $B=\op{diag}(\sigma_2,\sqrt{-1},-\sqrt{-1},-1,-1, \symbo, -\symbo)$, and
\item $C=\op{diag}(\sigma_2,-1,-1,\sqrt{-1},-\sqrt{-1}, \symbo, -\symbo)$
\end{itemize}
multiply to the identity and correspond to a ``non-abelian vertex'' of $P_8(3)$. These matrices are all conjugate to each other and correspond to the point $a=(\frac{1}{2}, \frac{1}{4},\frac{1}{4},\frac{1}{4},-\frac{1}{4},-\frac{1}{4},-\frac{1}{4}, -\frac{1}{2})$. The point $(a,a,a)\in P_8(3)$ comes from one of our basic extremal rays for the data
$d=2$, $\Gr(4,8)$, $I^1=\{2,3,4,7\}$, $I^2=I^3=\{1,3,4,7\}$ with $a=2$ and $j=1$.  We compute
$\langle\sigma_{I^1},\sigma_{I^2},\sigma_{I^3}\rangle_2=1$, and the level $\ell$ is computed as $4$, and
$$\lambda^1=\lambda^2=\lambda^3=(4,3,3,3,1,1,1,0)=\omega_1+2\omega_4+\omega_7.$$ It is easy to see that
$\frac{1}{\ell}\kappa(\lambda)=a$.
\subsection{The strange dual of Thaddeus' example}\label{E1}
\source{rigid-local-systems-multiplicative-eigenvalue-problem}
We compute the strange dual of the data consisting of $\lambda^1=\lambda^2=\lambda^3=(4,3,3,3,1,1,1,0)=\omega_1+2\omega_4+\omega_7$ at level $\ell=4$.
The data of $\mathcal{A}$ corresponding to the F-line bundle $\mathcal{B}(\vec{\lambda},\ell)$ on
$\Par_{8,\mathcal{O},S}$ (by Theorem \ref{egregium} and Proposition  \ref{corresp}) is $\mu^1=\mu^2=\mu^3=(7,4,4,1)$ with $n=8$.
This point corresponds to a rigid, irreducible unitary local system of rank $4$ on $\Bbb{P}^1-\{p_1,p_2,p_3\}$, with the same local monodromy ($\zeta^{-1},-1,-1,\zeta)$ at $p_1,p_2,p_3$, where
$\zeta=\zeta_8=\exp(2\pi\sqrt{-1}/8)$. By Prop \ref{propP}, this local system has infinite global monodromy.

There are 3 distinct eigenvalues at each point. The rigidity equation \eqref{egale} holds: $2 +(3-2)4^2= 3(2^2+1^2+1^2)$. This rigid local system belongs to the Goursat-III family of rigid local systems of rank $4$  (see Goursat's 1886 paper \cites{EGoursat}).

\iffalse
There are only two distinct Galois conjugates of our local system on  $\Bbb{P}^1-\{p_1,p_2,p_3\}$ corresponding to $\vec{a}$. Let $\vec{b}=(1/2,1/2,1/2,1/4,-1/4,-1/2,-1/2,-1/2)$. The non-trivial Galois conjugate is the point in $\Delta_8^3$ given by $(\vec{b},\vec{b},\vec{b})$. One can check that this is not a point of $P_9(3)$ by a direct computation. Hence the irreducible rigid local system corresponding to $\vec{a}$ has infinite global monodromy.
\fi


There is only one possible (optimal) Katz level lowering operation (see Section \ref{earlyfeedback}) for this example, because  at each point the eigenvalue with maximal multiplicity is unique. By a simple calculation, the resulting local system is of rank $2$, which is after a tensor product with a rank one local system, a local system on $\Bbb{P}^1-\{p_1,p_2,p_3\}$ with local monodromies $(-\zeta,-\zeta^{-1})=(\zeta^5,\zeta^4)$, $(1,\zeta^2)$ and $(1,\zeta^{-2})=(1,\zeta^6)$. This is a local system of hypergeometric type which is not unitary (see Theorem \ref{BHhere} below), since $(\zeta^{-5},\zeta^{-4})$ and
$(1,\zeta^2)$ do not interlace on the unit circle.
%Therefore the (Katz)  path down from this unitary local system is unique and passes through a non-unitary local system.

However  the choice of eigenvalues $-1$, $-1$ and $\zeta$ for the Katz operations at the three points gives rise to a rank $3$ unitary local system of hypergeometric type (after tensoring with a rank one)  with local monodromies $(-\zeta^{2},-1,1)=(\zeta^6,\zeta^4,1)$,$(-\zeta^{2},-1,1)=(\zeta^6,\zeta^4,1)$ and $(-\zeta^{-1},-\zeta^{-1},\zeta^{-2})=(\zeta^3,\zeta^3,\zeta^6)$. The unitarity is checked using Theorem \ref{BHhere}: Consider the local system corresponding to a twist by a rank one local system which has local monodromies $(1,1,\zeta^3)$, $(\zeta^6,\zeta^4,1)$ and $(\zeta,\zeta^3,\zeta^7)$. We may take $(\alpha_1,\alpha_2,\alpha_3)=(0,\frac{4}{8},\frac{6}{8})$ and $(\beta_1,\beta_2,\beta_3)=(1-\frac{7}{8},1-\frac{3}{8},1-\frac{1}{8})=(\frac{1}{8},\frac{5}{8},\frac{7}{8})$. The interlacing condition
therefore holds.  Also see Lemma \ref{hyperreduction}, and Question \ref{genphen}.
\iffalse
\subsection{An irreducible rigid local system of rank $5$}\label{E2}
\source{rigid-local-systems-multiplicative-eigenvalue-problem}
Consider the classical intersection number corresponding to $\Gr(5,8)$, $I^1=\{3,4,5,7,8\}$, $I^2=I^3=\{2,3,5,6,8\}$, $d=0$. It is easy to verify using the Littlewood-Richardson rule that this intersection number is one. So we are in a position to apply Theorems \ref{Adagio} and \ref{brahms} with $(a,j)=(3,1)$. Doing the enumerative computations of
Theorem \ref{brahms}, we find that $\ell=5$, $\lambda^1=\omega_2+2\omega_5 =(3,3,2,2,2,0,0,0)$ and $\lambda^2=\lambda^3=
2\omega_3+2\omega_6 =(4,4,4,2,2,2,0,0) $ (these are weights for $\op{SL}(8)$ at level $5$).

We now compute the corresponding unitary irreducible rigid local system: The corresponding data $\mathcal{A}$ is the following: $\mu^1=(5,5,2,0,0)$,
$\mu^2=\mu^3=(6,6,3,3,0)$ and $n=8$.  There are 3 distinct eigenvalues at each point with multiplicity decomposition $2,2,1$ (hence the local system is not  of hypergeometric or Pochhammer type). The rigidity equation $3(2^2+2^2+1^2)=(3-2)5^2+2$ (see equation \eqref{eqRi}) holds.
\fi
\subsection{An example that gives the determinant of cohomology  line bundle}
Consider  the case $d=0$, $1\in I^1$, $n\not\in I^1$. Consider $D(1,1)$, noting that  $d'=-1$.
The divisor $D(1,1)$ is clearly supported on  the divisor of $\Par_{n,\mathcal{O},S}$ where the underlying vector bundle is not trivial as a bundle (for otherwise it cannot have subsheaves of positive degree).  The corresponding line bundle is $\mathcal{D}$,  the determinant of cohomology. and hence $D(1,1)=\mathcal{D}$ a fact that follows from our numerical formulas as well.
It is interesting to note that we cannot have $1\in I^2, n\not\in I^2$  as well because the corresponding classical intersection number is zero since the line  $E^{p_1}_1$ is not  a subset of  $E^{p_2}_{n-1}$ for general flags.





\subsection{A rank $6$ rigid local system}\label{wilson}
\source{rigid-local-systems-multiplicative-eigenvalue-problem}
This example is due to A. Wilson. In the context of Proposition  \ref{practical}, consider the cycle $E=C(0,3,\mathcal{O},9,\vec{J})$ with $|S|=3$ on $\Par_{9,\mathcal{O},S}$ with
$J^1=\{2,6,9\}$, $J^2=J^3=\{3,6,9\}$. The line bundle $\mathcal{O}(E)=\mathcal{B}(\vec{\lambda},\ell)$ computed by Proposition  \ref{enumerative} has $\ell=6$ and $\lambda^1= 3\omega_2+2\omega_6$ and $\lambda^2=\lambda^3=2\omega_3+ 2\omega_6$. The inequality in Proposition  \ref{practical} (ii) holds and therefore the line bundle $\mathcal{B}(\vec{\lambda},\ell)$ is an F-line bundle on $\Par_{9,\mathcal{O},S}$. The corresponding unitary irreducible rigid local system is of rank $6$ with multiplicity decompositions $(3,2,1)$, $(2,2,2)$ and $(2,2,2)$ at $p_1,p_2$ and $p_3$. The rigidity equation \eqref{eqRi} holds:
$6^2+2=(3^2+2^2+1^2) + 2(2^2+2^2+2^2)$. The eigenvalues of the local monodromies at the 3 points are: at $p_1$, $\zeta^2$ (with multiplicity $3$) $\zeta^6$ (multiplicity $2$) and $1$ with multiplicity $1$; and at $p_2$ and $p_3$, the eigenvalues are $\zeta^3,\zeta^6$ and $1$ (with multiplicity $2$ each). Here $\zeta=\exp(2\pi\sqrt{-1}/9)$. The multiplicity of one at the first point shows that the cycle is of the form $\mathcal{O}(D(1,1))$ for $I^1=\{1,2,6\}$ and $I^2=I^3=\{3,6,9\}$, $d=1$ in Theorem \ref{Adagio} on $\Gr(3,9)$.

The Katz rank lowering algorithm (see Section \ref{earlyfeedback}) when applied to this rank six irreducible unitary rigid local system leads to irreducible rigid local systems of rank $5$ with non-semisimple local monodromy for some  of the choices in the algorithm (e.g., we pick the eigenvalues of maximal multiplicity to be $\zeta^2$, $1$ and $1$ at $p_1,p_2$ and $p_3$ respectively.

%There is a path down to rank ones in the Katz algorithm in which every local system that appears in the path has semi-simple local monodromy.
\subsection{An infinite collection of irreducible  unitary rigid local systems}\label{KO}
\source{rigid-local-systems-multiplicative-eigenvalue-problem}
J. Kiers and G. Orelowitz have reported the following series of examples that they have worked out in the context of  their joint work with S. Gao and A. Yong \cite{Kiers_etal}.

Let  $\Gr(k,n)$ with $n=3k-1$ for $k\ge2$, $I^1=I^3=\{2,5,8,\dots,3k-1\}$, $I^2=\{k,2k+1,2k+2,\dots 3k-1\}$ and $d=0$. The corresponding intersection number is computed to be one (this is a case of the Pieri formula).  In the setting of Theorems \ref{Adagio} and \ref{brahms}, consider $D(2k+1,2)=\mathcal{B}(\vec{\lambda},\ell)$. Here $d'=0$ and $J^1=J^3=I^1=I^3$, and $J^2$ is obtained by replacing the $2k+1$ in $I^2$ by $2k$.

After a careful combinatorial count of  intersection numbers (using the Littlewood-Richardson rule, the rim-hook formulas for the quantum numbers \cite{BCFF}, and Remark \ref{earlyAM}), Kiers and Orelowitz have obtained the following formulas: $\ell=(k-1)^2+1$,
$\lambda^1=\lambda^3= (k-1)(\omega_2 +\omega_5 + \omega_8+ \dots +\omega_{3k-4})$ and $\lambda^2=(k-1)\omega_{k} +\omega_{2k}$.

Setting $\zeta=\exp(2\pi\sqrt{-1}/n)$, the corresponding irreducible unitary rigid local system has rank $(k-1)^2+1$ and has local monodromy at $p_1$ and $p_3$ conjugate to the diagonal matrix with the following eigenvalues: $\zeta^2,\zeta^5,\zeta^8,\dots, \zeta^{3k-4}$ (each with multiplicity $k-1$) and $1$ (with multiplicity $1$). The local monodromy at $p_2$ is conjugate to the diagonal matrix with eigenvalues $\zeta^k$ (with multiplicity $k-1$), $\zeta^{2k}$ (with multiplicity $1$) and $1$ (with multiplicity $(k-1)(k-2)$).  The rigidity equation  \eqref{eqRi}  can be verified from these formulas.

For $k\ge 5$, these examples violate the speculation $\ell<n$ from an earlier version of this paper.





\subsection{An example of a vertex which is not an F-vertex}\label{Rex}
\source{rigid-local-systems-multiplicative-eigenvalue-problem}
The following example of Ressayre (reported in \cite[Section 10.5]{BHermit}) was made in the classical context:
Consider $G=\op{SL}(9)$, and  $\lambda=\omega_8+\omega_7+\omega_3$, $\mu=\nu=\omega_6+\omega_3$. The ray $\Bbb{Q}_{\ge 0}(\lambda,\mu,\nu)$ is an extremal ray of the tensor cone, but the rank of $(V_{\lambda}\tensor V_{\mu}\tensor V_{\nu})^G$ is $2>1$. When converted into a conformal block, the level is $4$, but the tuple $(\lambda,\mu,\nu,3)$ at level three has the same rank. Any level less than three will not
lead to a point in the fundamental alcove (for the first point) and therefore $(\lambda,\mu,\nu,3)$ gives an extremal ray of $\Par_{9,\mathcal{O,S}}$ with $|S|=3$ (use Remark \ref{classicall}, any expression of a multiple as a sum will need to project to the trivial decomposition in the classical context, so one of the summands should be a multiple of $\mathcal{D}$, making the other summand have levels that violate \eqref{weyl} in Theorem \ref{blacktea} at the first point),  and hence $(\lambda,\mu,\nu,3)$ leads to a vertex of $P_8(3)$.
% This vertex corresponds to  a family of reducible representations of the fundamental group but some of the factors are not rigid.
%Therefore not all vertices of $P_n(s)$ are rigid, but the building blocks have this property.

Analysing this example further as in \cite[Section 10.5]{BHermit}, we take the classical intersection $d=0$ and $I_1=\{3,7,8\}$, $I_2=I_3=\{3,6,9\}$ in $\Gr(3,9)$, and show that $(\lambda,\mu,\nu,3)$ is induced from the extremal ray corresponding to $\op{SL}(3)$, and the tuple $(\omega_2,\omega_2,\omega_2,1)$ (which is a basic extremal ray  coming from Theorem \ref{Adagio} for $\op{SL}(3)$).
